% =====================================================================
% Sección: Descripción y justificación del problema
% Responsable: Parte 1 — David Fernando Benjumea Mora
% =====================================================================
\section{Descripción y justificación del problema}
\label{sec:problema}

\subsection{Contexto}

Una red móvil ad hoc (MANET, \emph{Mobile Ad hoc Network}) es un conjunto de nodos
inalámbricos que se comunican entre sí sin depender de una infraestructura fija, como
puntos de acceso o estaciones base. Cada nodo actúa a la vez como terminal y como
enrutador: reenvía paquetes de otros nodos para que la información llegue a destinos
que están fuera de su alcance directo \cite{sarkar2013adhoc,zarate2021wireless}. Como
los nodos se mueven, la topología de la red cambia continuamente: los enlaces aparecen
y desaparecen, y las rutas deben descubrirse y repararse sobre la marcha.

Cuando el número de nodos crece, una organización plana deja de escalar: en protocolos
reactivos como AODV \cite{rfc3561aodv}, cada descubrimiento de ruta inunda la red con
solicitudes, y el tráfico de control termina compitiendo con el tráfico útil. La
alternativa habitual es una \textbf{arquitectura jerárquica}: los nodos se agrupan en
clústeres, cada clúster tiene un nodo coordinador o \emph{cluster head} (CH), y los CH
forman a su vez un nivel superior (\emph{backbone}) por el que viaja la comunicación
entre clústeres. Ospina y Ortiz \cite{ospina2017growth} estudiaron, con simulaciones en
ns-3, las condiciones bajo las cuales estas arquitecturas jerárquicas de dos y tres
capas son escalables, y mostraron que la relación de recursos entre capas sucesivas
puede describirse con un factor de crecimiento.

\subsection{El problema: dos niveles de movilidad}

En una MANET jerárquica la movilidad ocurre en dos niveles distintos:

\begin{itemize}
    \item \textbf{Movilidad de clúster:} el grupo completo se desplaza como un todo.
    Es el caso de un convoy de vehículos, una brigada de rescate o un pelotón que
    avanza junto.
    \item \textbf{Movilidad de nodo:} cada integrante se mueve individualmente dentro
    del área de su grupo, sin alejarse de él.
\end{itemize}

Los modelos de movilidad clásicos, como Random Waypoint o Random Walk, tratan a cada
nodo como independiente; los modelos de grupo, como RPGM (\emph{Reference Point Group
Mobility}), representan precisamente esta combinación de un movimiento colectivo con
desviaciones individuales \cite{camp2002survey,hong1999group}.

Ambos tipos de movilidad afectan a la red de forma diferente. La movilidad de nodo
cambia muchos enlaces, pero todos son \emph{locales}: rutas cortas dentro de un mismo
clúster, que AODV repara con poco costo. La movilidad de clúster cambia pocos enlaces,
pero son los enlaces \emph{críticos} entre los CH, por donde pasa todo el tráfico que
cruza la jerarquía; cuando uno de ellos se rompe, se interrumpen todos los flujos entre
clústeres que dependían de él. No es evidente cuál de los dos efectos domina, ni cómo
cambia esa relación al aumentar la velocidad.

El problema que aborda este trabajo es, entonces, \textbf{cuantificar y comparar, mediante
simulación en ns-3, el efecto de la movilidad de clúster y el de la movilidad de nodo
sobre el desempeño de una MANET jerárquica} de dos niveles (cuatro clústeres en el
primer nivel y dos superclústeres en el segundo), con enrutamiento AODV.

\subsection{Justificación}

\begin{itemize}
    \item \textbf{Diseño de redes jerárquicas.} Saber qué tipo de movilidad degrada más
    la red orienta las decisiones de diseño. Si domina la movilidad de clúster, conviene
    invertir en la estabilidad del backbone (elección de CH estables, gateways
    redundantes); si domina la de nodo, conviene invertir en el mantenimiento de rutas
    dentro de cada clúster.
    \item \textbf{Escenarios reales con movimiento en grupo.} Operaciones de rescate,
    redes militares tácticas, caravanas vehiculares y redes de sensores en animales
    presentan exactamente esta combinación de movimiento colectivo e individual.
    \item \textbf{Pertinencia de la simulación.} Montar un experimento de campo con
    decenas de equipos móviles es costoso y difícil de repetir. La simulación permite
    activar y desactivar cada tipo de movilidad por separado —algo imposible de controlar
    en el mundo real—, mantener constantes las demás variables y repetir cada
    experimento con distintas semillas aleatorias.
    \item \textbf{Relación con el curso.} La movilidad es un proceso estocástico, y cada
    corrida de la simulación es una realización de ese proceso. Por eso las conclusiones
    se obtienen con réplicas independientes, estimadores e intervalos de confianza, como
    se describe en la sección~\ref{sec:pregunta}.
\end{itemize}

\subsection{Alcance}

El estudio incluye una jerarquía fija de 20 nodos (cuatro clústeres de cinco nodos con
CH predefinido y dos superclústeres), el protocolo AODV, tráfico UDP de tasa constante
y cuatro escenarios de movilidad evaluados a cuatro velocidades. Quedan fuera del
alcance la reelección dinámica de CH, la comparación entre protocolos de enrutamiento,
el consumo de energía y el uso de la librería ns3-ai; estos aspectos se proponen como
trabajo futuro.
